\subsection{Dashboard}

La Dashboard è la prima schermata che l'utente vede dopo l'accesso e funge da pannello di controllo generale dell'esperienza su BrainBox. Nella versione precedente 
presentava statistiche aggregate, una sezione per gli esami in programma e un elenco delle attività recenti, ma con testi in parte in inglese e una struttura visiva poco organizzata. Nella versione riprogettata la Dashboard è stata tradotta interamente in italiano e arricchita: accoglie l'utente con un messaggio di benvenuto personalizzato, mostra una panoramica delle attività della settimana, gli esami imminenti con indicazione dei giorni rimanenti, le attività pianificate per oggi e gli appunti 
caricati di recente, tutti raggiungibili con un link diretto alla sezione corrispondente.

\paragraph{Design pattern applicati}

\begin{itemize}
    \item \textbf{Dashboard} (Tidwell): Il pattern prevede di raccogliere in un'unica 
    pagina i dati e le informazioni più rilevanti per l'utente, aggiornati in tempo reale, in modo che possa ottenere una panoramica immediata della propria situazione senza 
    dover navigare tra le diverse sezioni. Nella versione riprogettata la Dashboard centralizza in un colpo d'occhio lo stato degli appunti, degli esami imminenti, delle 
    attività di oggi e dell'attività settimanale, riducendo il bisogno di spostarsi tra le sezioni per raccogliere queste informazioni.

    \item \textbf{Titled Sections} (Tidwell): Il pattern prevede di suddividere una pagina densa di contenuti in sezioni distinte, ciascuna con un titolo visivamente riconoscibile, per facilitare la scansione e la comprensione immediata della struttura. Nella versione riprogettata la Dashboard è organizzata in blocchi chiaramente separati (\textit{Panoramica}, \textit{Questa settimana}, \textit{Esami}, 
    \textit{Attività di oggi} e \textit{Appunti recenti}) ognuno con titolo e contenuto omogeneo, rendendo immediato per l'utente identificare l'informazione che cerca.

    \item \textbf{One-Window Drilldown} (Tidwell): Il pattern prevede di mostrare una lista di elementi in una schermata e di permettere all'utente di accedere al dettaglio di ciascuno navigando nella sezione corrispondente. Nella Dashboard ogni sezione espone un riepilogo degli elementi più recenti o rilevanti (ultimi appunti, prossimi esami, attività di oggi) accompagnato da un link \textit{"Vedi tutti"} che porta direttamente alla sezione dedicata, evitando di sovraccaricare la Dashboard con informazioni complete.

\end{itemize}

\paragraph{Altre modifiche}
Il testo dell'interfaccia è stato interamente tradotto in italiano. Il pulsante 
\textit{"Passa a BrainBox Pro +"} (P17), presente nella versione precedente senza 
alcuna funzione attiva, è stato rimosso.

Problemi di usabilità risolti:
\begin{itemize}
  \item P11 (Mix italiano/inglese nella dashboard): Tutto il testo è stato tradotto 
  in italiano.
  \item P17 (Il pulsante "Passa a BrainBox Pro +" non è funzionante): Il pulsante è stato rimosso perché non era collegato a una funzione reale. Dato che la riprogettazione non ha riguardato il sistema degli abbonamenti, lasciarlo attivo avrebbe solo confuso l'utente con un comando inutile.
\end{itemize}